% from ega4-16-fr.tex
\oldpage[IV-4]{23}
déduits des homomorphismes canoniques
$P_{B/A}^n\to P_{S^{-1}B/A}^n$ (16.4.4), forment un système projectif et sont
bijectifs.
\end{proposition}

Il suffit de remarquer que
\[
  S^{-1}((B\otimes_A B)/\mathfrak{I}^{n+1})
  =S^{-1}(B\otimes_A B)/(S^{-1}\mathfrak{I})^{n+1}
\]
par platitude, et que
$S^{-1}(B\otimes_A B)=(S^{-1}B)\otimes_A(S^{-1}B)$ (I, 1.3.4).

\begin{corollary}[16.4.15]
\label{IV.16.4.15}
Les notations étant celles de (16.4.14), soit $R$ une partie multiplicative de
$A$ telle que $\rho(R)\subset S$. Alors on a des isomorphismes canoniques
\[
  S^{-1}P_{B/A}^n\xrightarrow{\sim}P_{S^{-1}B/R^{-1}A}^n
  \tag{16.4.15.1}
  \label{IV.16.4.15.1}
\]
formant un système projectif.
\end{corollary}

Il suffit évidemment de définir des isomorphismes canoniques
\[
  P_{S^{-1}B/A}^n\xrightarrow{\sim}P_{S^{-1}B/R^{-1}A}^n,
  \tag{16.4.15.2}
  \label{IV.16.4.15.2}
\]
c'est-à-dire que l'on est ramené au cas où $\rho(R)$ est formé d'éléments
\emph{inversibles} de $B$. Mais alors l'isomorphisme (16.4.15.2) est simplement
déduit de l'isomorphisme canonique
$B\otimes_A B\to B\otimes_{R^{-1}A}B$ par passage aux quotients
(0$_{\rm I}$, 1.5.3).

\begin{corollary}[16.4.16]
\label{IV.16.4.16}
Soient $f:X\to S$ un morphisme de préschémas, $x$ un point de $X$,
$s=f(x)$. Alors on a des isomorphismes canoniques
\[
  (\mathscr{P}_{X/S}^n)_x\xrightarrow{\sim}
  P_{\mathcal{O}_x/\mathcal{O}_s}^n
  \tag{16.4.16.1}
  \label{IV.16.4.16.1}
\]
formant un système projectif.
\end{corollary}

On déduit de là des isomorphismes pour les gradués associés, et en particulier
un isomorphisme canonique
\[
  (\Omega_{X/S}^1)_x\xrightarrow{\sim}
  \Omega_{\mathcal{O}_x/\mathcal{O}_s}^1.
  \tag{16.4.16.2}
  \label{IV.16.4.16.2}
\]

\begin{corollary}[16.4.17]
\label{IV.16.4.17}
Soient $k$ un corps, $K$ le corps des fractions rationnelles
$k(T_1,\ldots,T_r)$. Alors, pour tout entier $n$, l'homomorphisme de
$K[U_1,\ldots,U_r]$ ($U_i$ indéterminées) dans $P_{K/k}^n$ qui, à tout $U_i$,
fait correspondre $d^nT_i-T_i\cdot1$, est surjectif et définit un isomorphisme
du quotient $K[U_1,\ldots,U_r]/\mathfrak{m}^{n+1}$ (où $\mathfrak{m}$ est
l'idéal engendré par les $U_i$) sur $P_{K/k}^n$.
\end{corollary}

Cela résulte de (16.4.8), (16.4.10) et (16.4.14), où l'on fait $A=k$,
$B=k[T_1,\ldots,T_r]$ et $S=B-\{0\}$.

On retrouve ainsi le fait que les $dT_i$ forment une base du $K$-espace
vectoriel $\Omega_{K/k}^1$ (0, 20.5.10).

\begin{proposition}[16.4.18]
\label{IV.16.4.18}
Soient $f:X\to Y$, $g:Y\to Z$ deux morphismes de préschémas; et considérons
les homomorphismes canoniques de $\mathcal{O}_X$-Algèbres augmentées
(16.4.3.3)
\[
  g_{X/Y/Z}:\mathscr{P}_{X/Z}^n\longrightarrow\mathscr{P}_{X/Y}^n
  \tag{16.4.18.1}
  \label{IV.16.4.18.1}
\]
\[
  f_{X/Y/Z}:f^*(\mathscr{P}_{Y/Z}^n)
  \longrightarrow\mathscr{P}_{X/Z}^n.
  \tag{16.4.18.2}
  \label{IV.16.4.18.2}
\]
Alors $g_{X/Y/Z}$ est surjectif, et son noyau est l'Idéal engendré par l'image
par $f_{X/Y/Z}$ de l'Idéal d'augmentation de
$f^*(\mathscr{P}_{Y/Z}^n)$.
\end{proposition}

